\begin{proofbox}
	\textbf{\textcolor{CProof}{思路提示}}

	利用 $A+B = \pi - C$ 以及 $\frac{A+B}{2} = \frac{\pi}{2} - \frac{C}{2}$，套用诱导公式即可证得。

	\medskip
	\textbf{\textcolor{CProof}{正式推导}}
	\begin{description}[style=nextline, leftmargin=2.8em, labelsep=0.5em, itemsep=0.55em]
		\item[\textbf{步骤一（补角代换）：}]
			\[
				A+B = \pi - C
			\]
			\par\textbf{理由：} 由内角和条件 $A+B+C=\pi$ 移项即得。

		\item[\textbf{步骤二（正弦公式）：}]
			\[
				\begin{aligned}
					\sin(A+B) &= \sin(\pi - C) \\
					&= \sin C
			\end{aligned}
			\]
			\par\textbf{理由：} 诱导公式 $\sin(\pi - \alpha)=\sin\alpha$。

		\item[\textbf{步骤三（余弦公式）：}]
			\[
				\begin{aligned}
					\cos(A+B) &= \cos(\pi - C) \\
					&= -\cos C
			\end{aligned}
			\]
			\par\textbf{理由：} 诱导公式 $\cos(\pi - \alpha)=-\cos\alpha$。

		\item[\textbf{步骤四（正切公式）：}]
			\[
				\begin{aligned}
					\tan(A+B) &= \tan(\pi - C) \\
					&= -\tan C \quad (C\neq\frac{\pi}{2})
			\end{aligned}
			\]
			\par\textbf{理由：} 诱导公式 $\tan(\pi - \alpha)=-\tan\alpha$；需保证 $C\neq\frac{\pi}{2}$ 使分母非零。

		\item[\textbf{步骤五（半角转换）：}]
			\[
				\begin{aligned}
					\frac{A+B}{2} &= \frac{\pi - C}{2} \\
					&= \frac{\pi}{2} - \frac{C}{2}
			\end{aligned}
			\]
			\par\textbf{理由：} 将 $A+B = \pi - C$ 两边同除以 $2$。

		\item[\textbf{步骤六（半角正弦）：}]
			\[
				\begin{aligned}
					\sin\frac{A+B}{2} &= \sin\left(\frac{\pi}{2} - \frac{C}{2}\right) \\
					&= \cos\frac{C}{2}
			\end{aligned}
			\]
			\par\textbf{理由：} 诱导公式 $\sin(\frac{\pi}{2} - \beta)=\cos\beta$。

		\item[\textbf{步骤七（半角余弦）：}]
			\[
				\begin{aligned}
					\cos\frac{A+B}{2} &= \cos\left(\frac{\pi}{2} - \frac{C}{2}\right) \\
					&= \sin\frac{C}{2}
			\end{aligned}
			\]
			\par\textbf{理由：} 诱导公式 $\cos(\frac{\pi}{2} - \beta)=\sin\beta$。
	\end{description}

	\medskip
	\textbf{\textcolor{CProof}{结论回扣}}

	\textbf{由以上七步，完整证明了三角形内角和诱导恒等变换的四条核心公式（含正切条件）以及两条半角公式。}
\end{proofbox}
